\section{Appendix}\label{sec:appendix}


\subsection{Functional analysis and the FEM}\label{appendix:functional_analysis_and_the_fem}
\paragraph{Lebesgue space.} Let $\Omega \subset \mathbb{R}^d$ be an open set and $\mathbb{F}$ denote either $\mathbb{R}$ or $\mathbb{C}$. The space $L^2(\Omega)$, often referred to as the Lebesgue space of square integrable functions, is defined as 
\begin{equation*}
    L^2(\Omega) =
    \left\{
        f : \Omega \to \mathbb{F} \, \Bigg| \,
        \int_{\Omega} |f(x)|^2 \, dx < \infty
    \right\},
\end{equation*}
with inner product
\begin{equation*}
    \langle f,\, g \rangle_{L^2(\Omega)} = \int_\Omega f(x) \overline{g(x)} \, dx.
\end{equation*}
The space $L^2(\Omega)$ is a Hilbert space, i.e., a complete inner product space. For details, see Section~3 in Kreyszig \cite{kreyszigIntroductoryFunctionalAnalysis1989}.


\paragraph{Sobolev space.} The space $H^1(\Omega)$, known as a first-order Sobolev space, is defined as
\begin{equation*}
    H^1(\Omega) =
    \left\{
        f \in L^2(\Omega) \, \Bigg| \,
        \frac{\partial f}{\partial x_i} \in L^2(\Omega), \, i = 1, \dots, d  
    \right\},
\end{equation*}
with inner product
\begin{equation*}
    \langle f,\, g \rangle_{H^1(\Omega)} = 
    \langle f,\, g \rangle_{L^2(\Omega)} +
    \langle \nabla f,\, \nabla g \rangle_{L^2(\Omega)}.
\end{equation*}
The space $H^1(\Omega)$ is a Hilbert space and a subspace of $L^2(\Omega)$. See Section~2 in Quarteroni for further details \cite{quarteroniNumericalModelsDifferential2014}.


\paragraph{The stiffness and the mass matrix.}\label{appendix:stiffness}%
Let $V_h \subset H^1(\Omega)$ be the finite element space defined in Equation~\eqref{eq:fem_space}. Under the basis $\{\phi_i\}_{i=1}^N$, the stiffness matrix $A_h \in \mathbb{R}^{N \times N}$ and mass matrix $M \in \mathbb{R}^{N \times N}$ are defined as
\begin{equation*}
    (A_h)_{ij} = \int_\Omega 
    \left( \sigma \nabla \phi_j \cdot \nabla \phi_i 
    + c \phi_j \phi_i \right) dx,
    \quad
    M_{ij} = \int_\Omega \phi_j \phi_i \, dx.
\end{equation*}
Both matrices are symmetric positive definite (Appendix~\ref{appendix:linear_algebra}) under the standing assumptions that $\sigma$ is positive definite and $c > 0$, see Theorem~3.4 and Theorem~4.10 in Larson and Bengzon \cite{larsonFiniteElementMethod2013}. The boundary mass matrix $M_\partial \in \mathbb{R}^{N_b \times N_b}$ is given by
\begin{equation*}
    (M_\partial)_{ij} = \int_{\partial\Omega} \phi_i \phi_j \, ds,
\end{equation*}
where the indices $i, j$ tabulate over the boundary indices of the mesh. Like the mass matrix $M$, the boundary mass matrix $M_\partial$ is SPD \cite{larsonFiniteElementMethod2013}.

$A_h$, $M$ and $M_\partial$ are sparse matrices, since $(A_h)_{ij}$, $M_{ij}$ and $(M_\partial)_{ij}$ are non-zero only if $\phi_i$ and $\phi_j$ share an element in the mesh.

The assembly of stiffness matrices and mass matrices in two dimensions can be done through methods of numerical integration, as described in Section~3.5 in Larson and Bengzon \cite{larsonFiniteElementMethod2013}.


\subsection{Derivation of the adjoint formulation}\label{appendix:adjoint_derivation}
Consider the forward operator $\mathcal{K} = \mathcal{T} \circ \mathcal{S}$ given by Equation~\eqref{eq:K}, and the variational formulation of $\mathcal{S}$ by Equation~\eqref{eq:S_cont}. Define the bilinear form
\begin{equation*}
    a(u,\, v) =
    \int_\Omega (\sigma \nabla u \cdot \nabla v + c u v) \, dx,
\end{equation*}
then, the variational formulation of $\mathcal{S}$ reads as follows: find $u \in H^1(\Omega)$ such that
\begin{equation}\label{eq:bilinear_form}
    a(u,\, v) = \int_\Omega f v \, dx.
\end{equation}

For some $g \in L^2(\partial\Omega)$, we seek $w = \mathcal{K}^* g$. By the definition of $\mathcal{K}^*$,
\begin{equation*}
    \langle \mathcal{K} f,\, g \rangle_{L^2(\partial\Omega)} =
    \langle f,\, \mathcal{K}^* g \rangle_{L^2(\Omega)}.
\end{equation*}
and since $\mathcal{K} f = u|_{\partial\Omega}$, we have that
\begin{equation}\label{eq:adjoint_identity_1}
    \langle \mathcal{K} f,\, g \rangle_{L^2(\partial\Omega)} =
    \int_{\partial\Omega} u g \, ds.
\end{equation}
In Equation~\eqref{eq:bilinear_form}, choose $v = w = \mathcal{K}^* g$. We see that
\begin{equation}\label{eq:adjoint_identity_2}
    a(u, w) = \int_{\Omega} f w \, dx = \langle f,\, \mathcal{K}^* g \rangle_{L^2(\Omega)}.
\end{equation}
Inserting Equations~\eqref{eq:adjoint_identity_1} and~\eqref{eq:adjoint_identity_2} into the definition of $\mathcal{K}^*$, and using the identity $a^*(v,\, u) = a(u,\, v)$ (by the definition of $a^*$ \cite{kreyszigIntroductoryFunctionalAnalysis1989}),
\begin{equation}\label{eq:adjoint_identity_3}
    a^*(w,\, u) = \int_{\partial\Omega} g u \, ds.
\end{equation}
Since Equation~\eqref{eq:adjoint_identity_3} must hold for all $u \in H^1(\Omega)$, it defines a variational problem for $w$. And since $a$ is symmetric ($\sigma$ is symmetric),
$a^*(w,\, u) = a(u,\, w) = a(w,\, u)$, we get that Equation~\eqref{eq:adjoint_identity_3} reduces to:
find $w \in H^1(\Omega)$ such that
\begin{equation*}
    a(w,\, u) = \int_{\partial\Omega} g u \, ds
    \quad \text{for all } u \in H^1(\Omega),
\end{equation*}
where
\begin{equation*}
    a(w, u) = \int_{\Omega}
    (\sigma \nabla w \cdot \nabla u + c w u)\, dx.
\end{equation*}
This is the variational formulation of Equation~\eqref{eq:variational_formulation_adjoint} with $u := v$.


\subsection{Linear algebra and matrix decompositions}\label{appendix:linear_algebra}
The definitions and properties throughout this subsection follow Golub and Van Loan \cite{golubMatrixComputations4th2013}.

\paragraph{Symmetric positive definiteness.}
Let $A \in \mathbb{R}^{n \times n}$. $A$ is said to be symmetric if $A = A^T$. If $A$ is symmetric and if 
\begin{equation*}
    x^T A x > 0
\end{equation*}
for all nonzero $x \in \mathbb{R}^n$, then $A$ is said to be symmetric positive definite.


\paragraph{QR factorization.}
The QR factorization of $A\in\mathbb R^{m\times n}$ is a decomposition
\begin{equation*}
    A = QR,
\end{equation*}
where $Q\in\mathbb R^{m\times n}$ has orthonormal columns and $R\in\mathbb R^{n\times n}$ is upper triangular. The matrix $Q$ forms an orthonormal basis for $\operatorname{range}(A)$.


\paragraph{LU decomposition.} The LU decomposition of $A \in \mathbb{R}^{n \times n}$ is a factorization
\begin{equation*}
    A = L U,
\end{equation*}
where $L$ is a lower triangular matrix and $U$ an upper triangular matrix. $A$ has an LU decomposition under the condition that $\operatorname{det}(A[1:k, 1:k]) \neq 0$ for $k=1,\dots, n-1$. More generally, any square matrix $A$ has an LU decomposition with partial pivoting (LUP), $PA = LU$. Computing the LU decomposition costs $\mathcal{O}(n^3)$ flops, specifically, $\tfrac{2}{3}n^3$.


\paragraph{Cholesky factorization.} Assume that $A\in\mathbb R^{n\times n}$ is SPD. Then $A$ has a unique Cholesky factorization
\begin{equation*}
    A = L L^T,
\end{equation*}
where $L$ is a lower triangular matrix with positive diagonal elements. This is a special case of the LU decomposition, with $U = L^T$. The cost of this procedure is $\mathcal{O}(n^3)$ flops, specifically, $\tfrac{1}{3}n^3$, making it faster than standard LU decomposition.


\paragraph{Incomplete Cholesky factorization.} If $A \in \mathbb{R}^{n\times n}$ is SPD with Cholesky factorization $A = L L^T$, then the incomplete Cholesky factorization is an approximation
\begin{equation*}
    A \approx \tilde L \tilde L^T,
\end{equation*}
where $\tilde L$ is a lower triangular matrix that is sparser than $L$. There are many strategies for computing the incomplete Cholesky, see Lin and Moré \cite{linIncompleteCholeskyFactorizations1999}. The main advantage over the full Cholesky factorization is reduced storage and computational cost, making it particularly effective for large sparse SPD matrices.


\subsection{Solving triangular systems}\label{appendix:triangular}
Let $R \in \mathbb{R}^{n \times n}$ be triangular. Solving $R x = y$ is done by forward (or backward) substitution: starting from the first (or last) row, each $x_i$ is solved directly after substituting all previously computed values. If $R$ is sparse, zero entries are skipped, and by Lemma~2.2.1 in George and Liu \cite{georgeComputerSolutionLarge1984} the cost reduces to $\mathcal{O}(\operatorname{nnz}(R))$, where $\operatorname{nnz}(R)$ is the number of nonzeros in $R$.


\paragraph{General linear systems.} Let $A \in \mathbb{R}^{n\times n}$ be invertible and consider $A x = y$, then the solution is $x = A^{-1} y$. In practice, the standard procedure is to compute the LU factorization $A = LU$, reducing the problem to two triangular systems
\begin{equation*}
    L z = y, \quad U x = z \implies x = U^{-1} z = U^{-1} L^{-1} y.
\end{equation*}
Each triangular system can be solved by forward and backward substitution at $\mathcal{O}(n^2)$ \cite{golubMatrixComputations4th2013}. The factorization itself costs $\tfrac{2}{3}n^3$ flops, so the problem scales as $\mathcal{O}(n^3)$. If $A$ is symmetric positive definite, $U = L^T$ and the Cholesky factorization $A = LL^T$ halves the cost to $\tfrac{1}{3}n^3$ (Appendix~\ref{appendix:linear_algebra}).


\paragraph{Sparse linear systems.} For sparse systems $A x = y$, computing the LU factorization comes with additional challenges. Matrix entries that are zero in $A$ but non-zero in $L$ or $U$, so-called fill-ins, require more memory and result in a more expensive factorization \cite{larsonFiniteElementMethod2013}. Additionally, the cost of solving the resulting triangular systems is higher. Much research has gone into the development of algorithms which avoid fill-ins, see Davis \cite{davisDirectMethodsSparse2006} for an overview.


\paragraph{The nested dissection method.} For sparse SPD matrices of size $N \times N$ arising from 2D elliptic PDEs (mass and stiffness matrices), the structure of the matrices can be exploited. Alan George proposed the nested dissection method in 1973 \cite{georgeNestedDissectionRegular1973}, a divide and conquer approach that computes Cholesky factorization in $\mathcal{O}(N^{3/2})$ flops. The resulting triangular matrix $L$ contains $\mathcal{O}(N \log N)$ non-zero elements, so solving $L x = y$ costs $\mathcal{O}(\operatorname{nnz}(L)) = \mathcal{O}(N \log N)$.


\subsection{Dynamical low-rank approximation}\label{appendix:dynamical_low_rank_approximation}
In \cite{kochDynamicalLowrankApproximation2007}, Koch and Lubich study the problem of finding a rank-$r$ approximation $X(t)$ to a time-dependent data matrix $A(t)$,
\begin{equation*}
    \min_{X(t) \in \mathcal{M}_r} \| X(t) - A(t) \|_F,
\end{equation*}
where $\mathcal{M}_r = \{M \in \mathbb{R}^{m \times n} : \operatorname{rank}(M) = r \}$. Instead of computing a best approximation via SVD at every time $t$, they propose a dynamical approach. They seek a low-rank matrix $Y(t) \in \mathcal{M}_r$ such that its derivative $\dot{Y}(t)$ approximates $\dot{A}(t)$:
\begin{equation}\label{eq:dlra1}
    \min_{\dot{Y}(t) \in \mathcal{T}_{Y(t)}(\mathcal{M}_r)} \| \dot{Y}(t) - \dot{A}(t) \|_F,
\end{equation}
where $\mathcal{T}_{Y(t)}(\mathcal{M}_r)$ denotes the tangent space of $\mathcal{M}_r$ at $Y(t)$.

Let $Y = U \Sigma V^T$ be a factorization where $U \in \mathbb{R}^{m \times r}$ and $V \in \mathbb{R}^{n \times r}$ have orthonormal columns, and $\Sigma \in \mathbb{R}^{r \times r}$ is non-singular (using the SVD yields a diagonal $\Sigma$, but this is not a necessary assumption). Koch and Lubich show that Equation~\eqref{eq:dlra1} is equivalent to the Galerkin condition: find $\dot{Y} \in \mathcal{T}_{Y}(\mathcal{M}_r)$ such that
\begin{equation*}
    \langle \dot{Y} - \dot{A},\, \delta Y \rangle_F = 0,\quad \text{for all } \delta Y \in \mathcal{T}_{Y}(\mathcal{M}_r).
\end{equation*}
By defining orthogonal projectors $P_U^{\perp} = I - U U^T$ and $P_V^{\perp} = I - V V^T$ they prove (Proposition~2.1) that this condition yields the following system of matrix differential equations:
\begin{equation*}
    \begin{aligned}
    \dot{\Sigma} &= U^T \dot{A} V, \\
    \dot{U} &= P_U^{\perp} \dot{A} V \Sigma^{-1}, \\
    \dot{V} &= P_V^{\perp} \dot{A}^T U \Sigma^{-T}.
    \end{aligned}
\end{equation*}
These equations allow for evolution of the low-rank factors using numerical integration techniques.


\subsection{Conjugate gradient algorithm}\label{appendix:conjugate_gradient_algorithm}
Consider the quadratic function $f(x) = \tfrac 1 2 x^T A x - y^T x$ where $A$ is SPD, which has the gradient $\nabla f(x) = A x - y$. By the optimality condition $\nabla f(x) = 0$, the minimum of the function is the solution of $A x = y$.

Consider the update rule $x_{k+1} = x_k + \alpha_k d_k$, where $\alpha_k$ is a step-size and $d_k$ some search direction. The optimal $\alpha_k$ is the minimizer of the function $h(\alpha) = f(x_{k} + \alpha d_k)$, and setting $h'(\alpha) = 0$ with $d_k = -\nabla f(x_k)$ yields the step-size
\begin{equation*}
    \alpha_k = \frac{r_k^Tr_k}{d_k^TAd_k}, \quad r_k = y - Ax_k.
\end{equation*}
This is the steepest descent method \cite{chong2023introduction}. However, if the condition number of $A$ is high, then the contours of $f$ get stretched out, giving an elliptic landscape. Using the steepest descent rule, we have that $d_{k+1}^T d_k = 0$ \cite{chong2023introduction}, so the search directions will zig-zag across the landscape.

Instead, if we seek search directions that are $A$-orthogonal, $d_{k+1}^T A d_k = 0$, we are moving within a transformed landscape where the contours of $f$ are perfectly circular. Choosing $d_{k+1}$ as a linear combination of $-\nabla f(x_{k+1}) = r_{k+1}$ and the previous direction $d_k$,
\begin{equation*}
    d_{k+1} = r_{k+1} + \beta_k d_k,
\end{equation*}
enforcing $d_{k+1}^T A d_k = 0$,
\begin{equation*}
    (r_{k+1} + \beta_k d_k)^T A d_k = 0,
\end{equation*}
and solving for $\beta_k$, leads to
\begin{equation*}
    \beta_k = -
    \frac{r_{k+1}^T A d_k}
	     {d_k^T A d_k}.
\end{equation*}
This is the conjugate gradient algorithm \cite{chong2023introduction}.

The $A$-term in the formula of $\beta_k$ can be computationally expensive to compute if $f$ is non-linear \cite{chong2023introduction}. The formula of $\beta_k$ can be rewritten without $A$ to obtain the Fletcher-Reeves formula \cite{fletcherFunctionMinimizationConjugate1964},
\begin{equation*}
    \beta_k =
    \frac{r_{k+1}^T r_{k+1}}
	     {r_k^T r_k}.
\end{equation*}
Alternatively, the Hestenes-Stiefel formula \cite{hestenesMethodsConjugateGradients1952} computes $\beta_k$ as
\begin{equation*}
    \beta_k = -
    \frac{r_{k+1}^T (r_{k+1} - r_k)}
	     {d_k^T(r_{k+1} - r_k)}.
\end{equation*}
Another well-known modification is the Polak-Ribière formula \cite{polakNoteConvergenceMethodes1969},
\begin{equation*}
    \beta_k = 
    \frac{r_{k+1}^T(r_{k+1} - r_k)}
         {r_k^T r_k}.
\end{equation*}


\subsection{Performance metrics}\label{appendix:performance_metrics}
Let $x \in \mathbb{R}^N$ denote the true source (its coefficient vector) and $\hat x \in \mathbb{R}^N$ its corresponding reconstruction. To evaluate reconstruction quality, we consider three performance metrics: Euclidean error, EMD, and AUC-IoU.


\paragraph{Earth mover's distance.} The earth mover's distance measures the dissimilarity between two distributions, which in our case refers to $x$ and $\hat x$. Let $x_d$ and $\hat x_d$ be the normalized distributions of $x$ and $\hat x$, respectively, computed as
\begin{equation*}
    (x_d)_i = \frac{|x_i|}{\sum_{j=1}^N |x_j|}, \quad
    (\hat x_d)_i = \frac{|\hat x_i|}{\sum_{j=1}^N |\hat x_j|}.
\end{equation*}
Let $p_i \in \mathbb{R}^2$ denote the spatial coordinates of the $i$-th node in the mesh. The EMD between $x_d$ and $\hat x_d$ is
\begin{equation*}
    \operatorname{EMD}(x_d, \hat x_d) =
    \inf_{\gamma \in \Pi}
    \sum_{i,j=1}^N \gamma_{ij}\, \| p_i - p_j \|,
\end{equation*}
where $\Pi$ is the set of all joint distributions with marginals $x_d$ and $\hat x_d$ \cite{rubnerMetricDistributionsApplications1998}.


\paragraph{AUC-IoU.} The IoU, known as the Jaccard index in statistics, measures the similarity between two sets. It is defined as the size of intersection divided by the size of the union \cite{raschkaPythonMachineLearning}.

We define the binary mask of $x$ at threshold $\tau \in (0, 1)$ as the set
\begin{equation*}
    S_{\tau}(x) = \{ i : x_i \ge \tau \operatorname{max}(x) \}.
\end{equation*}
Then, the IoU of $x$ and $\hat x$ at threshold $\tau$ is the IoU of their binary masks,
\begin{equation*}
    \operatorname{IoU}(\tau) =
    \frac{|S_\tau(x) \cap S_\tau(\hat x)|}{|S_\tau(x) \cup S_\tau(\hat x)|}.
\end{equation*}
Here, $| \cdot |$ denotes the number of elements of a set (its cardinality).

The AUC-IoU is then the integration of the $\operatorname{IoU}(\tau)$ curve,
\begin{equation*}
    \operatorname{AUC-IoU} = \int_{0}^{1} \operatorname{IoU}(\tau) \, d\tau.
\end{equation*}
This can be approximated using the Trapezoidal rule. The AUC-IoU lies in $[0, 1]$, with 1 being a perfect score.